% Reserved D.69 candidate. Unregistered pending independent proof audit.
\subsection{Generic rigidity in the full profinite cusp packet}

\begin{statement}{D.69}{Proved}{Generic completed cusp classes have singleton full profinite fibres}
Let \(F=\pi_1((\mathbb P^1\setminus\{0,1,\infty\})_{\overline{\mathbb Q}})\)
be the actual free profinite geometric group, with the
unit tangential arithmetic action at zero. Put
\[
 \widehat{\mathcal K}=H^1(G_{\mathbb Q},\widehat{\mathbb Z}(1))
          =\prod_\ell\mathcal K_\ell,
 \qquad \mathcal K_\ell=H^1(G_{\mathbb Q},\mathbb Z_\ell(1)).
                                                        \tag{D.69.1}
\]
For odd \(\ell\), give \(\mathcal K_\ell\) its
uniform \(\ell\)-adic \(c_0\) topology; at two use
\[
 \mathcal K_2=\mathbb Z/2\ \oplus\
                   c_0(\{q\text{ rational prime}\},\mathbb Z_2),
                                                        \tag{D.69.2}
\]
with the sign factor discrete and the other factor
uniform two-adic. Give (D.69.1) the product topology.
There is a dense \(G_\delta\) subset
\(\widehat{\mathcal G}\subset\widehat{\mathcal K}\)
such that, for each \(a\in\widehat{\mathcal G}\),
the actual full profinite section fibre
\[
 H^1_{\rm cts}(G_{\mathbb Q},F)
       \longrightarrow H^1(G_{\mathbb Q},F^{\rm ab})
                      =\widehat{\mathcal K}^{2}
       \qquad\text{over }(a,0)
                                                        \tag{D.69.3}
\]
is the singleton consisting of its completed cusp
section at zero. The analogous assertion holds in
each of the other two rational cusp packets.

The parameters may be chosen infinitely ramified
in every prime component. At every prime \(\ell\),
the corresponding full geometric pro-\(\ell\)
section fibre is itself a singleton. At five the
set may simultaneously satisfy D.66--D.68, including
the singleton metabelian fibre and the non-effective
finite-stage obstruction away from its cusp.
No classification of all abelian pairs or of all
full profinite sections is asserted.
\end{statement}

\begin{proof}
We first check the coefficient topology and the
prime two; neither is inferred from the five-adic
statement by an implicit convention.

\paragraph{Power completion and the sign cocycle.}
The group \(\mathbb Q^\times\) is the product of
its sign group with the free abelian group on
the positive rational primes. Finite Kummer theory,
followed by the inverse limit of the finite
coefficient groups, gives
\[
 \mathcal K_\ell\simeq
 \begin{cases}
 c_0(\{q\},\mathbb Z_\ell),&\ell\ne2,\\
 \mathbb Z/2\oplus c_0(\{q\},\mathbb Z_2),&\ell=2.
 \end{cases}                                           \tag{D.69.4}
\]
Here the finite Kummer quotients at each exponent
are discrete direct sums over the rational primes;
their inverse limit is \(c_0\), not the unrestricted
product over those primes. The invariant groups
of the finite root-of-unity coefficients are
finite and Mittag--Leffler, so that continuous
\(H^1\) agrees with this inverse limit.
Also continuous cochains with values in
\(\widehat{\mathbb Z}(1)=\prod_\ell\mathbb Z_\ell(1)\)
are exactly the products of the component cochains.
Products are exact for these abelian complexes,
which proves (D.69.1). Equivalently it is the
power completion
\(\varprojlim_m\mathbb Q^\times/(\mathbb Q^\times)^m\)
with its inverse-limit topology, by the Chinese
remainder theorem. These Kummer groups need not
be compact: the finite-exponent quotients are
infinite discrete groups.

Fix compatible roots of every positive prime,
as in D.66, at each coefficient prime. At two
fix compatible primitive roots of unity. A
continuous representative for the sign class is
\[
            \kappa(-1)(g)=\frac{\chi_2(g)-1}{2}
                          \quad\text{in }\mathbb Z_2(1).
                                                        \tag{D.69.5}
\]
Indeed take \(\zeta_{2^{n+1}}\) as a compatible
\(2^n\)-th root of \(-1\). Its Kummer exponent
relative to \(\zeta_{2^n}\) is exactly (D.69.5)
modulo \(2^n\). Twice this cocycle is the
coboundary \(\chi_2-1\). Thus (D.69.5) represents
the nonzero class of order two; no linear section
of that order-two cohomology group into cocycles
is claimed. On each fixed-sign clopen component
we use the representative
\[
             a(g)=\varepsilon\kappa(-1)(g)
                         +\sum_q a_q\kappa(q)(g),
                    \qquad \varepsilon\in\{0,1\}.
                                                        \tag{D.69.6}
\]
The sum is uniformly convergent. Deep changes in
the prime coefficients leave the chosen sign
representative unchanged. For odd \(\ell\),
use the same formula without its sign term.

Every \(\mathcal K_\ell\) is a complete
separable metric group: the uniform \(c_0\)
argument of D.66 applies for every prime, and
at two there are just two clopen copies of it.
Their countable product is complete and
separable in a usual bounded product metric.
The Baire argument below therefore applies in
the precise topology specified in the statement.

\paragraph{The nonabelian construction at an arbitrary prime.}
Fix \(\ell\), let \(P_\ell=F^{(\ell)}\),
and put \(H_\ell=P_\ell'\). The completed
cusp parameter gives the honest action
\[
       \sigma_a(g)=\operatorname{Int}(x^{a(g)})\sigma_0(g),
       \qquad\sigma_0(g)(x)=x^{\chi_\ell(g)}.
                                                        \tag{D.69.7}
\]
As in D.67, the exact fibre above \((a,0)\)
is \(H^1(G_{\mathbb Q},H_\ell;\sigma_a)\).
The conjugator comparison uses
\(H^0(G_{\mathbb Q},\mathbb Z_\ell(1)^2)=0\),
which holds at two as well.

At every rational prime \(q\ne\ell\), the
tripod's marked smooth model and unit tangent
have good reduction. The specialization square
of \cite[Definition A.3 and Remark A.3.1]{DE_Hoshi2012}
is cartesian for \(\Sigma=\{\ell\}\), since the
residue characteristic is excluded from that
set. Proposition A.6 of the same primary source
then gives the specialized inertia alternative
used in D.67. Its hypotheses allow \(\ell=2\)
and every \(q\ne\ell\).

The Magnus proof of
\(C_{P_\ell}(x^b)=\langle x\rangle\) for
\(b\ne0\) works over \(\mathbb Q_\ell\) for
every prime. Abelianization turns the normalizer
of any nonzero subgroup of this cusp into that
centralizer. Thus the full local section is in
a cusp packet whenever its abelian inertia is
nonzero. For the pair \((a,0)\), that cusp is
zero. Equality of the actual abelian cocycles
puts the local conjugator in \(H_\ell\): the
local fixed abelian image is zero because
Frobenius acts by \(q\), with \(q-1\ne0\)
in the torsion-free lattice \(\mathbb Z_\ell^2\).
Consequently every integral nonabelian cocycle
\(c:G_{\mathbb Q}\to H_\ell\) satisfies
\[
 \begin{cases}
 c|_{G_{\mathbb Q_q}}\text{ is an integral }H_\ell
                         \text{-coboundary},& a_q\ne0,\\
 c(I_q)=1,&a_q=0.
 \end{cases}                                           \tag{D.69.8}
\]
The sign term of (D.69.6) is unramified at odd
\(q\), as are all local unit Kummer classes
at coefficient prime different from the residue
characteristic. Thus it changes no inertia
coordinate in (D.69.8).

Use the iterated Frattini quotients
\(Q_r=P_\ell/\Phi^r(P_\ell)\), of exponent
dividing \(\ell^r\), and the actual image
\(H_r\) of \(H_\ell\) therein. With
\[
 S_r(a)=\{\ell,\infty\}\cup
              \{q:a_q\notin\ell^r\mathbb Z_\ell\},
                                                        \tag{D.69.9}
\]
every reduced integral cocycle lies in the
finite actual set
\(Z^1(G_{\mathbb Q,S_r(a)},H_r;\rho_{a,r})\).
The proof is exactly the integral assertion
(D.69.8): at a later nonzero inertia prime,
its coboundary is for an inner power divisible
by \(\ell^r\), hence trivial on \(Q_r\).
At zero inertia it is already trivial.
The finite sets and actions are constant under
changes by \(\ell^r c_0\), on a fixed sign
component when \(\ell=2\).

For every reduced cocycle nontrivial on the
action kernel, apply Chebotarev to its actual
finite affine image and choose a prime with
identity action and nonidentity cocycle value.
Avoid the finite support of the low-valuation
coefficients of the original parameter, and
add \(\ell^M\kappa(q)\), with \(M\ge r\)
and arbitrarily large. This creates a nonzero
inertia coordinate of valuation exactly \(M\)
without changing any of the finite data.
The finite affine image really contains the
chosen kernel element, as in D.67; hence
\cite[Chapter V, Theorem 3.23]{PBC_MilneCFT}
applies without an assumption about the existence
of a Frobenius class. The finitely many tests
at each precision define an open dense set,
including in each fixed-sign component at two.

Intersect these sets over \(r\), and also
require nonzero inertia away from \(\ell\).
For every surviving parameter all cocycles
vanish on the finite action kernels, by
(D.69.8) at their witness primes. They descend
continuously to the actual compact action
image on \(H_\ell\). A nonzero inertia prime
gives a normal subgroup of inner powers of
\(x^{a_q}\), whose fixed subgroup in
\(H_\ell\) is trivial. Its restriction is
an integral coboundary by (D.69.8); conjugate
by its element of \(H_\ell\) to make that
restriction trivial. Comparing the cocycle at
\(ng=g(g^{-1}ng)\) makes every remaining
value fixed by that normal subgroup, hence
trivial. This proves a dense \(G_\delta\)
set \(\mathcal G_\ell\subset\mathcal K_\ell\)
with singleton full pro-\(\ell\) fibre.
Intersect it with the complement of the
countable union of finite-support submodules.
It remains dense \(G_\delta\), and every
parameter is infinitely ramified. At two
take the union of the constructions on the
two clopen sign components. At five we may
also intersect with the loci already proved
in D.66--D.68.

\paragraph{The full profinite comparison is a theorem about sections.}
Put
\[
                   \widehat{\mathcal G}
                           =\prod_\ell\mathcal G_\ell.
                                                        \tag{D.69.10}
\]
This is dense \(G_\delta\) in (D.69.1).
Indeed it is a countable intersection of
inverse images of dense open sets under the
open coordinate projections; every basic
product neighborhood restricts only finitely
many coefficient primes. Its prescribed cusp
section exists in the full profinite group:
the product of the component Kummer cocycles
is a continuous \(\widehat{\mathbb Z}(1)\)-cocycle,
which is inserted into the actual inertia
subgroup \(\langle x\rangle\) of the rational cusp.

Let \(s\) be any full profinite section with
abelian class \((a,0)\), for
\(a\in\widehat{\mathcal G}\). Its image in
every actual geometrically pro-\(\ell\)
quotient has the same class as this fixed
full cusp section, by the singleton result
just proved. D.24 applies to precisely these
two actual full sections: one is arbitrary,
and the other is a fixed member of the
packet of a rational cusp. It concludes
that their full section classes agree.
Its required local point-or-cusp theorem is
the full-formation case of D.56, as already
recorded in D.24. This is not an assertion
that the geometric pro-prime quotients
jointly detect every element of \(F\).
It uses the proved primewise-to-full theorem
for section classes and gives exactly the
singleton in (D.69.3).

Finally, the \(\mathbb Q\)-defined permutations
of the three punctures transport the actual
arithmetic section sets and cusp packets.
On packet parameters they give continuous
affine bijections, including the possible
sign constant at two. They therefore carry
the dense \(G_\delta\) just constructed to
the corresponding loci for the other two
cusps. No assertion outside these packet
abelian fibres is used or deduced.
\end{proof}
